\documentclass[10pt,letterpaper]{article}
\usepackage{cornell-notes}

\title{Chapter 28: Labeled Trees}
\author{}
\date{}

\setDocTitle{Chapter 28: Labeled Trees}
\setDocSubtitle{Combinatorial Algorithms}
\setDocVersion{v1.0}
\setDocStatus{Study Companion}
\setDocOwner{Combinatorial Algorithms Cornell Notes}
\setDocAuthor{}
\setDocDate{\today}
\setDocSummary{Cornell-format study and review notes for Chapter 28: Labeled Trees.}

\setCornellCollection{Combinatorial Algorithms Cornell Notes}
\setCornellUnitType{Chapter}
\setCornellUnitNumber{28}
\setCornellUnitTitle{Labeled Trees}
\setCornellCompanionLabel{Study and review companion}

\hypersetup{pdftitle={Chapter 28: Labeled Trees},pdfsubject={Cornell notes},pdfauthor={}}

\begin{document}
\maketitle
\begin{tcolorbox}[colback=Paper,colframe=Ink]{\small\bfseries\color{Teal}CORNELL NOTES}\par{\LARGE\bfseries Chapter 28: Labeled Trees (LBLTRE)}\par\medskip\textbf{Topic:} Pr\"ufer coding, enumeration, and uniform generation\hfill\textbf{Date:} \today\par\textbf{Study purpose:} Use a bijective code to represent and generate labeled trees.\end{tcolorbox}
\begin{tcolorbox}[colback=Teal!5,colframe=Teal,title=Essential question]Why does deleting the smallest leaf create a reversible length-$n-2$ code, and what does that code reveal about labeled trees?\end{tcolorbox}
\begin{tcolorbox}[colback=white,colframe=Mist,title=Learning objectives]\begin{itemize}\item Encode and decode a labeled tree.\item Derive Cayley's count.\item Read vertex degrees from code multiplicities.\item Explain the in-place linear-time implementation.\item Generate a uniform random labeled tree.\end{itemize}\textbf{Prerequisites:} trees, leaves, bijections, elementary probability.\end{tcolorbox}
\section*{Cornell notes}
\begin{longtable}{@{}Q!{\color{Mist}\vrule width 1pt}A@{}}\cornellhead
\cue{What is the code?}{A tree on labels $1,\ldots,n$ corresponds to a word $(p_1,\ldots,p_{n-2})$ over those labels. The encoding repeatedly removes the smallest currently available leaf and records its unique neighbor.}
\cue{How is it decoded?}{At step $r$, choose the smallest label not appearing in the unprocessed suffix $p_r,\ldots,p_{n-2}$, join it to $p_r$, and delete that leaf and code symbol. After all symbols are consumed, join the two remaining labels.}
\cue{Why is decoding valid?}{A label absent from the remaining code must be a leaf of the not-yet-built portion. Choosing the smallest such label reverses the encoder's deterministic tie rule. Induction shows that every word yields one tree and re-encoding returns the same word.}
\cue{Cayley's formula}{There are $n$ choices for each of $n-2$ code positions. The bijection therefore proves that the number of labeled trees on $n$ vertices is $n^{,n-2}$.}
\cue{Degrees from the code}{Vertex $v$ appears in the code exactly $\deg(v)-1$ times. Hence $\deg(v)=1+\#\{r:p_r=v\}$. Leaves are precisely the labels absent from the code, and prescribed degree counts follow from multinomial word counts.}
\cue{Uniform random generation}{Choose the $n-2$ code symbols independently and uniformly from $\{1,\ldots,n\}$, then decode. Because the mapping is bijective, every labeled tree receives probability $n^{-(n-2)}$.}
\cue{Useful consequence}{The probability that a fixed vertex is a leaf is $(1-1/n)^{n-2}$. Therefore the expected number of leaves is $n(1-1/n)^{n-2}$, asymptotic to $n/e$.}
\cue{LBLTRE state scheme}{The in-place decoder marks last occurrences in the code and uses TREE states: a negative marker means the label still occurs later, $0$ means it is currently eligible as the smallest leaf, and a positive entry records its selected neighbor/edge. Moving pointers find the next eligible label without rescanning all suffixes.}
\cue{Complexity and checks}{Each code position and label changes state only a constant number of times, giving $O(n)$ time and $O(n)$ storage. Validate $n\geq2$ and every symbol in $1..n$; the final output must have exactly $n-1$ edges and be connected.}
\cue{Retrieval practice}{\begin{enumerate}\item Which leaf is deleted during encoding?\item How is the next leaf found in decoding?\item How many times does a degree-$d$ vertex occur?\item Why is a random code uniform over trees?\item What makes LBLTRE linear-time?\end{enumerate}}
\end{longtable}
\begin{tcolorbox}[colback=Ink!4,colframe=Ink,title=Cornell summary]
Pr\"ufer coding is a bijection between labeled trees and length-$n-2$ words. Encoding deletes the smallest leaf and records its neighbor; decoding repeatedly reconnects the smallest label absent from the remaining suffix. The bijection proves $n^{n-2}$ trees, turns degrees into symbol multiplicities, and makes uniform random generation immediate. LBLTRE implements the inverse mapping in linear time by marking last occurrences and reusing an array to distinguish labels that still occur, labels eligible as leaves, and completed parent links.
\end{tcolorbox}
\end{document}
